%% LyX 2.4.4 created this file.  For more info, see https://www.lyx.org/.
%% Do not edit unless you really know what you are doing.
\documentclass[english]{article}
\usepackage[T1]{fontenc}
\usepackage[latin9]{inputenc}
\usepackage{enumitem}
\usepackage{amsmath}
\usepackage{amsthm}
\usepackage{geometry}
\geometry{verbose,tmargin=1in,bmargin=1in,lmargin=1in,rmargin=1in}

\makeatletter
%%%%%%%%%%%%%%%%%%%%%%%%%%%%%% Textclass specific LaTeX commands.
\numberwithin{equation}{section}
\numberwithin{figure}{section}
\newlength{\lyxlabelwidth}      % auxiliary length 

%%%%%%%%%%%%%%%%%%%%%%%%%%%%%% User specified LaTeX commands.
\usepackage{fancyhdr}
\usepackage{lastpage}
\pagestyle{fancy}
\usepackage{enumitem}
\usepackage{ifthen}
\everymath=\expandafter{\the\everymath\displaystyle}
%\DeclareMathSizes{10}{10}{10}{10}
\usepackage{multicol}

\usepackage{tikz}
\usetikzlibrary{patterns}
\usetikzlibrary{plotmarks}
\usepackage{pgfplots}
\pgfplotsset{compat=newest}

\usepackage{calc}
\usepackage[nomessages]{fp}

\usepackage{datetime}
% Added by lyx2lyx
\setlength{\parskip}{\smallskipamount}
\setlength{\parindent}{0pt}

\makeatother

\usepackage{babel}
\begin{document}
\renewcommand{\author}{Dr.\,James Ashe}

\newcommand{\classNum}{232}

\newcommand{\classSec}{A and B}

\newcommand{\examType}{Exam 3}

\renewcommand{\date}{Due: \formatdate{10}{4}{2013}}

%%First page's header%%%%%%%%%%%%%%%%%%%%%%
\fancypagestyle{firststyle} %%header settings for first pagr
{
	\fancyhead[L]{MTH \classNum{}(\classSec{}) \author{}\\
	\textbf{\examType{}} ~\date{}}
	\fancyhead[C]{}
	\fancyhead[R]{\textbf{Name:\hspace{4cm}}}
	\setlength{\headsep}{14pt}
}
%%%%%%%%%%%%%%%%%%%%%%%%%%%%%%%%%%%%%%%%%%

%%Next page's header%%%%%%%%%%%%%%%%%%%%%%%
\setlist{nolistsep}
\lhead{\classNum{}(\classSec{})} 
\chead{\textbf{Name:}} 
\cfoot{\thepage{}~of~\pageref{LastPage}} 
\setlength{\headsep}{5pt} 
%%%%%%%%%%%%%%%%%%%%%%%%%%%%%%%%%%%%%%%%%%%

%%Various settings; see the preamble for other settings%%%%%
%\setenumerate[0]{leftmargin=12pt,itemindent=0pt,labelwidth=0pt}
\setlist{topsep=.5pt}
\renewcommand{\labelenumi}{\textbf{\arabic{enumi}.}}
\renewcommand{\labelenumii}{\textbf{\alph{enumii}.}}
\thispagestyle{firststyle}
%%%%%%%%%%%%%%%%%%%%%%%%%%%%%%%%%%%%%%%%%%

\newcommand{\xmax}{0}
\newcommand{\xmin}{-\xmax}
\newcommand{\xstep}{0}
\newcommand{\ymax}{0}
\newcommand{\ymin}{-\ymax}
\newcommand{\ystep}{0} 

\textbf{Justify your answers and show your work. }Unless indicated
otherwise, each problem is worth 5 points.

\begin{multicols}{2}
\begin{enumerate}[resume]
\item Let $f(x)=\sqrt{x+5}$ for $x\geq-5$ {[}3 pts each{]}

\begin{enumerate}
\item Find the domain of $f^{-1}$.\vfill
\item Find the range of $f^{-1}$.\vfill\columnbreak
\item Graph $f^{-1}$.\\
\renewcommand{\xmax}{10} \renewcommand{\xstep}{5}
\renewcommand{\ymax}{\xmax} \renewcommand{\ystep}{5}
%%%%%%%
\begin{tikzpicture} 
	\begin{axis}[
		axis x line=middle, axis y line=middle,     		
		%xtick={-\xmax,-\xstep,...,\xmax},
		%ytick={-\ymax,-\ystep,...,\ymax},
		minor x tick num={4},
		minor y tick num={4},
		grid=none, %both
		xlabel={$$},     
		ylabel={$$},
		xmin=-\xmax.25,xmax=\xmax.25,
		ymin=-\ymax.25,ymax=\ymax.25,     		
		width=5.5cm,
		height=5.5cm] 		
		\addplot[domain=-5:9.5,thick,samples=1000,draw=red,->]{sqrt(x+5)} node[above] {$f$};
		%\addplot[domain=\xmin:1,thick,samples=100,draw=red]{-x-4} node[right] {$$};
		%\addplot[color=red,mark=o,solid] coordinates {(1,3)};
		%\addplot[color=red,mark=*,solid] coordinates {(1,-5)};
		%\addplot[mark=noner,smooth,domain=-1:1]{}; 
	\end{axis}
\end{tikzpicture} 
\end{enumerate}
\end{enumerate}
\end{multicols}
\begin{enumerate}[resume]
\item Using the graph of $f(x)$ shown below to complete the following.

\begin{enumerate}
\item {[}4 pts{]} Find two intervals on which $f$ has an inverse, making
each interval as large as possible.\\
\renewcommand{\xmax}{4.7123889}
\renewcommand{\xmin}{-.25}
\renewcommand{\xstep}{0} %must be xmin+n
\renewcommand{\ymax}{1.25}
\renewcommand{\ymin}{-1.35}
\renewcommand{\ystep}{0} %must be ymin+n
%%%%%%%
\begin{tikzpicture}
\begin{axis}[height=4cm, 
	width=5cm, 
	axis lines=middle,
	xtick={-6.28318, -4.7123889, -3.14159, -1.5708,1.5708, 	3.14159, 4.7123889, 6.28318},
	xticklabels={$-2\pi$, $-\frac{3\pi}{2}$, $-\pi$, 	$\frac{\pi}{2}$,        $\frac{\pi}{2}$, $\pi$, 	$\frac{3\pi}{2}$, $2\pi$},
	%axis line style={-latex}, 
	xmin=\xmin,xmax=\xmax, 
	ymin=\ymin,ymax=\ymax,
	%restrict y to domain=-1:10,
	no markers,
	xlabel={$x$},ylabel={$y$}, 
	%every axis x label/.append style={anchor=north}, 
	%every axis y label/.append style={anchor=east}, 
	area style,tension=0.08]
		
	\addplot[domain=0:2*pi,thick,smooth,red,samples=100] {cos(deg(x))}; 
	
	%\node at (axis cs:5,7) {$f$};
\end{axis}
\end{tikzpicture}
\item {[}2 pts{]} Assuming that $f^{-1}$ exists on one of these intervals,
find $f^{-1}\left(1\right)$.\vspace{1cm}
\end{enumerate}
\item Given the function $f(x)=\left(x-2\right)^{2}$ for $x\geq2$, find
the slope of the line tangent to the graph of $f^{-1}$ at the point
$\left(1,3\right)$.\vfill
\end{enumerate}
\emph{Perform the indicated operation.}
\begin{enumerate}[resume]
\item $\frac{d}{dx}\left(e^{-3x}+5e^{x}\right)$\vfill
\item $\frac{d}{dx}$$x^{2}\ln x-\pi$\vfill
\end{enumerate}
\newpage\emph{Perform the indicated operation.}
\begin{enumerate}[resume]
\item ${\displaystyle \int\left(10e^{5x}+\frac{1}{x}\right)\,dx}$\vfill
\item ${\displaystyle \int\frac{dx}{3x-5}}$\vfill
\item Use logarithmic differentiation to find the derivative of $f(x)=\left(1+e^{x}\right)^{x}$.
\vfill
\item Solve the equation for $x$.

\begin{enumerate}[resume]
\item {[}3 pts{]} $e^{2x}=30$\vfill
\item {[}3 pts{]} $\ln2x=5$\vfill
\end{enumerate}
\item $\frac{d}{dx}\log_{3}x$\vfill
\item $\frac{d}{dx}\frac{3^{x}}{x^{2}}$\vfill
\end{enumerate}
\newpage\emph{Perform the indicated operation.}
\begin{enumerate}[resume]
\item $\frac{d}{dx}\left(x^{\pi}+1\right)$\vfill\vspace{-3cm}
\item {[}8 pts{]}${\displaystyle \int}2xe^{2x}\,dx$\vfill
\item {[}8 pts{]}${\displaystyle \int x\cos\left(8x\right)\,dx}$\vfill
\item {[}8 pts{]}${\displaystyle \int10x\ln(2x)\,dx}$\vfill
\end{enumerate}
\emph{Bonus:}
\begin{enumerate}[resume]
\item ${\displaystyle \int x\cdot\sin x\cos x\,dx}$
\end{enumerate}

\end{document}
